\section{Nonreversible conservative motion of cuts}
\label{sec:nonreversible}

The reduction in Theorem~\ref{thm:spatial-gap} extends to conservative
jumps that permute the sites, without requiring inverse jumps to have
the same rate.  The state in this section is the full nonempty cut set.
As in the spatial model, unordered interval lengths need not form a
Markov process.

For $\sigma\in S_n$ put $(T_\sigma F)(S)=F(\sigma S)$, and let
\begin{equation}\label{eq:permutation-motion}
 H=\sum_{\sigma\in S_n}r_\sigma(T_\sigma-I),
 \qquad r_\sigma\ge0,\qquad
 L=L_B+DH,\quad D\ge0.
\end{equation}
Terms with $r_\sigma=0$ may be omitted.  No relation between
$r_\sigma$ and $r_{\sigma^{-1}}$ is imposed.
The induced one-coordinate generator is
\begin{equation}\label{eq:one-coordinate-motion}
 (H_1a)(x)=\sum_\sigma r_\sigma\bigl(a(\sigma x)-a(x)\bigr).
\end{equation}
It preserves the subspace
$\mathbb C^n_0=\{a:\sum_xa(x)=0\}$.
Adjoints below are with respect to the uniform measure on coordinates
or $\pi_B$ on cut sets, as appropriate.  Define
\begin{align}
 \kappa(H_1)
 &=\min\{\operatorname{Re}z:
             z\in\operatorname{Spec}(-H_1|_{\mathbb C^n_0})\},\label{eq:kappa-real}\\
 \kappa_{\mathrm{s}}(H_1)
 &=\lambda_{\min}\!\left(
       -\tfrac12(H_1+H_1^*)|_{\mathbb C^n_0}\right).\label{eq:kappa-symmetric}
\end{align}
Zeros are included in these minima.  Thus the definitions also apply
when the coordinate motion alone has several communicating classes.

\begin{theorem}\label{thm:nonreversible}
For every $n\ge2$, $\beta>0$, $D\ge0$, and choice of the rates in
\eqref{eq:permutation-motion}, the law $\pi_B$ in
\eqref{proof:eq:piB} is invariant for $L$.  On the mean-zero subspace,
the asymptotic spectral rate is exactly
\begin{equation}\label{eq:nonreversible-real-gap}
 g_{\mathrm{as}}
 :=\min\{\operatorname{Re}z:
            z\in\operatorname{Spec}(-L),\ z\ne0\}
 =\min\{\gamma_Q,\ \eta_1+D\kappa(H_1)\}.
\end{equation}
The largest constant $g_{\mathrm{s}}$ for which
\begin{equation}\label{eq:nonreversible-contraction}
 \|e^{tL}F\|_{L^2(\pi_B)}
 \le e^{-g_{\mathrm{s}}t}\|F\|_{L^2(\pi_B)}
 \qquad(t\ge0,\ \mathbb E_{\pi_B}F=0)
\end{equation}
is
\begin{equation}\label{eq:nonreversible-symmetric-gap}
 g_{\mathrm{s}}=\min\{\gamma_Q,\ \eta_1+D\kappa_{\mathrm{s}}(H_1)\}.
\end{equation}
Here $\eta_1$ is the smallest eigenvalue of $-\mathcal L_1$, as
in Theorem~\ref{thm:spatial-gap}.
\end{theorem}

\begin{proof}
Every coordinate permutation preserves $\pi_B$, which is constant on
each cardinality layer.  Hence $H$ preserves $\pi_B$, and
$T_\sigma^*=T_{\sigma^{-1}}$.  The reaction generator is irreducible,
so adding the nonnegative jump rates $DH$ preserves irreducibility.
In particular, zero is the simple stationary eigenvalue of $L$.

The Boolean harmonic decomposition is equivariant under every
coordinate permutation.  On degree $\ell$, the reaction generator
acts on the radial factor, and $H$ acts on the harmonic factor.
Thus, with $H_\ell$ denoting the latter action,
\begin{equation}\label{eq:nonreversible-tensor}
 L|_\ell=\mathcal L_\ell\otimes I+I\otimes DH_\ell.
\end{equation}
The inner product factors into the radial reversible inner product
and the harmonic inner product.  Each permutation is unitary on the
harmonic factor, so for every complex vector $v$ in that factor,
\[
 \operatorname{Re}\langle v,-H_\ell v\rangle
 =\frac12\sum_\sigma r_\sigma
       \|T_\sigma v-v\|^2\ge0.
\]
Consequently every eigenvalue of $-H_\ell$ has nonnegative real part.

Diagonalizing the self-adjoint radial factor in
\eqref{eq:nonreversible-tensor} expresses its spectrum as the sums of
radial eigenvalues and $D$ times eigenvalues of $-H_\ell$.
This statement does not require diagonalizability of $H_\ell$.
In degree zero, $H_0=0$, and the first positive eigenvalue is
$\gamma_Q$.  By Proposition~\ref{prop:higher-exact}, all reaction
eigenvalues in degrees $\ell\ge2$ are at least $\gamma_Q$.
These sectors therefore cannot lower the minimum real part.

In degree one, $H_\ell$ is precisely the coordinate operator
\eqref{eq:one-coordinate-motion} on $\mathbb C^n_0$:
if $F_a(S)=\sum_{x\in S}a(x)$, then $HF_a=F_{H_1a}$.
The smallest radial eigenvalue is $\eta_1$.
Both terms on the right of \eqref{eq:nonreversible-real-gap} are
attained by eigenvalues, proving that formula.

For the contraction statement, the symmetric part is
$L_{\mathrm{s}}=L_B+D(H+H^*)/2$.
It is a reversible permutation-jump generator plus $L_B$.
The same sector argument gives its spectral gap as
\eqref{eq:nonreversible-symmetric-gap}.  For
$F_t=e^{tL}F$,
\[
 \frac{d}{dt}\|F_t\|_2^2
 =2\operatorname{Re}\langle F_t,LF_t\rangle
 =2\langle F_t,L_{\mathrm{s}}F_t\rangle
 \le-2g_{\mathrm{s}}\|F_t\|_2^2.
\]
Integration proves \eqref{eq:nonreversible-contraction}.
Conversely, differentiating any proposed prefactor-one bound at
$t=0$ gives the Poincar\'e inequality for $-L_{\mathrm{s}}$.
Its gap is therefore the largest possible such constant.
\end{proof}

\begin{remark}[Asymptotic rate and finite-time norm bounds]
In finite dimension,
\[
 \lim_{t\to\infty}-\frac1t
       \log\|e^{tL}|_{L^2_0(\pi_B)}\|=g_{\mathrm{as}}.
\]
For every $\varepsilon>0$ there is a finite $C_\varepsilon$ such that
the norm is at most
$C_\varepsilon e^{-(g_{\mathrm{as}}-\varepsilon)t}$.
If a defective eigenvalue lies on the line with real part
$-g_{\mathrm{as}}$, a polynomial factor in $t$ can occur.
Thus \eqref{eq:nonreversible-real-gap} alone does not assert a
prefactor-one estimate at that exponent.  Formula
\eqref{eq:nonreversible-symmetric-gap} identifies the exact exponent
for that stronger type of estimate.
\end{remark}

\subsection{An explicit strict increase of the asymptotic rate}

On three coordinates, take
\[
 H=\tfrac12(T_{(12)}-I)+\tfrac12(T_{(23)}-I)
       +(T_{(123)}-I).
\]
The convention $(123):1\mapsto2\mapsto3\mapsto1$ gives
\begin{equation}\label{eq:three-site-matrix}
 H_1=
 \begin{pmatrix}
  -3/2&3/2&0\\
  1/2&-2&3/2\\
  1&1/2&-3/2
 \end{pmatrix}.
\end{equation}
Its rows and columns sum to zero, but it is not symmetric.
The characteristic polynomials of $-H_1$ and its symmetric part are
\[
 z\bigl(z^2-5z+27/4\bigr),
 \qquad z(z-2)(z-3),
\]
respectively.  Hence
\[
 \kappa(H_1)=\frac52,\qquad
 \kappa_{\mathrm{s}}(H_1)=2,
 \qquad
 \operatorname{Spec}(-H_1)
   =\{0,\,5/2+i/\sqrt2,\,5/2-i/\sqrt2\}.
\]
The two rates in Theorem~\ref{thm:nonreversible} become
\[
 g_{\mathrm{as}}=\min\{\gamma_Q,\eta_1+5D/2\},
 \qquad
 g_{\mathrm{s}}=\min\{\gamma_Q,\eta_1+2D\}.
\]
For $0<D<2(\gamma_Q-\eta_1)/5$, the asymptotic rate is strictly
larger than the reversible symmetric-part rate.

The same example extends to arbitrarily many sites.
Let $V=\{1,2,3\}\times W$, where $W$ carries a connected symmetric
random-walk generator $G$ with gap at least $5/2$.
At each $w\in W$, apply the three permutations above inside
$\{1,2,3\}\times\{w\}$ with the displayed rates.
For every graph edge $wz$ and every $i\in\{1,2,3\}$, also transpose
$(i,w)$ and $(i,z)$ at that edge's rate.
The one-coordinate generator is then
\[
 H_1^{V}=H_1\otimes I+I\otimes G.
\]
Its eigenvalues add, so its two relevant rates remain $5/2$ and $2$.
The preceding exact formulas therefore hold for every
$n=3|W|$, with $\gamma_Q$ and $\eta_1$ evaluated at that $n$.
All conservative moves in this construction affect two or three sites.

\begin{remark}[Scope of the conservative extension]
The rates in \eqref{eq:permutation-motion} specify unconditional
permutation jumps; an unchanged occupancy configuration simply
contributes zero.  Occupancy-dependent directed exclusion rates are
not covered by this argument.  In that model, the action on a
linear coordinate observable generally contains quadratic occupancy
terms, and the degree-one tensor reduction used here need not hold.
\end{remark}
